\section{Method}
\subsection{Formulation}
\paragraph{Conditional Flow matching for audio-driven video generation}
Flow matching video generative models \cite{flowmatching, chen2025hunyuanvideo, wan} adopts a neural network to generate realistic video frames by modeling a timestep-dependent vector field that transports samples from a noise distribution to a target video distribution. Given a ground truth conditional video distribution $q(\boldsymbol{x} \vert \boldsymbol{c})$ where $\boldsymbol{x} \in \mathbb{R}^{t \times h \times w \times c}$ is the encoded video latent. $\boldsymbol{c} = \{y, a, \boldsymbol{x}_{\mathrm{ref}}\}, y \in \mathbb{R}^{m \times d_{\mathrm{text}}}, a \in \mathbb{R}^{n \times d_{\mathrm{audio}}}, \boldsymbol{x}_{\mathrm{ref}} \in \mathbb{R}^{t_{\mathrm{ref}} \times h\times w \times c}$ are the conditions, including the text prompt embedding, and the audio embedding, and the reference frames latent. Conditional flow matching defines a series of distributions by interpolating $q(\boldsymbol{x} \vert y, a)$ with a known trivial distribution (e.g. Gaussian noise) $p(\boldsymbol{x} \vert y, a)$ using a continuous variable $t \in [0. 1]$.

\begin{equation}
q_t(\boldsymbol{x} \vert y, a) = (1 - t) \cdot p(\boldsymbol{x} \vert y, a) + t \cdot q(\boldsymbol{x} \vert y, a).
\end{equation}

To be specific, a random variable $\boldsymbol{x}_t \sim q_t(\boldsymbol{x} \vert y, a)$ can be obtained by interpolating between $\boldsymbol{x}_0 \sim p(\boldsymbol{x} \vert y, a)$ and $\boldsymbol{x}_1 \sim q(\boldsymbol{x} \vert y, a)$ via $\boldsymbol{x}_t = (1 - t) \cdot \boldsymbol{x}_1 + t \cdot \boldsymbol{x}_0$. The generative model $\boldsymbol{v}_\theta(\cdot)$, parameterized by $\theta$, is trained to match a continuous velocity field $\boldsymbol{v}_\theta(\boldsymbol{x}_t \vert y, a) \simeq \frac{d \boldsymbol{x}_t}{d t}$. To achieve so, we adopt the conditional flow matching objective. 

\begin{equation}
    \mathcal{L}_{\mathrm{fm}} = \mathbb{E}_{t, \boldsymbol{x}_0, \boldsymbol{x}_1} \|\boldsymbol{v}_\theta(\boldsymbol{x}_t \vert y) - (\boldsymbol{x}_1 - \boldsymbol{x}_0)\|_2^2.
\end{equation}

 An ODE solver can be used to sample from a flow matching generative models.
% Preliminary and formulation: A brief introduction on streaming audio-driven image-to-video flow matching generative model. Will need to include the contents:
% (1) Mathmatical formation and brief introduction of flow matching model.
% (2) Input and output formulation, how the streaming is achieved. (In our model, we use audio $a$ (using wav2vec embedding), image $x_{ref}$, streaming context frames $x_{context}$ as condition to produce new video chunks).

\paragraph{Sparse-frame video dubbing}
% Given a source video latent $\boldsymbol{x}_0 \in \mathbb{R}^{t \times h \times w \times c}$ and target audio $a \in \mathbb{R}^{n \times d_{\mathrm{audio }}}$, sparse-frame video dubbing aims to achieve audio-driven video-to-video editing. During the process, the model references only the selected keyframe $\boldsymbol{x}_{\mathrm{ref}}$ to keep identity, key gestures, and camera movement of the source video, while giving the model freedom to generate audio-aligned lip, head, and body movements. Crucially, this paradigm operates on infinite long sequences, necessitating generative capabilities beyond short clips to maintain synchronization across extended durations.
Video dubbing localizes content by replacing original audio with translated speech while preserving visual authenticity. As formalized in this work, the task transforms a source video latent $\boldsymbol{x}_0 \in \mathbb{R}^{t \times h \times w \times c}$ and a target audio $a \in \mathbb{R}^{n \times d_{\mathrm{audio }}}$ into an output video where lip movements, facial expressions, and body dynamics synchronize organically with the new audio.
Traditional video dubbing techniques focus exclusively on oral region inpainting—editing lip movements while freezing head rotations, facial expressions, and body gestures \cite{li2024latentsync}. This creates immersion-breaking mismatches, as static body language contradicts emotional speech (e.g., a rigid posture during passionate dialogue). Sparse-frame video dubbing, illustrated in \cref{fig:teaser}, fundamentally redefines this process: it preserves only select keyframes $\boldsymbol{x}_{\mathrm{ref}}$ to anchor identity, emotional cadence, symbolic gestures, and camera trajectories—critical for visual continuity—while liberating full-body dynamics (facial expressions, head motions, body gestures) to organically synchronize with dubbed audio. As \cref{fig:teaser} demonstrates, this paradigm shift enables lifelike alignment where head turns follow speech rhythm and gestures amplify emotional tone—impossible with lip-only editing. Crucially, sparse-frame dubbing operates on infinite-length sequences, demanding generative continuation beyond short clips to maintain synchronization across extended durations, a capability unattainable with traditional frame-by-frame inpainting.

% Recent advances in audio-driven generative models have significantly improved lip synchronization capabilities for video dubbing \cite{li2024latentsync}. However, these methods predominantly focus on oral region inpainting, resulting in mismatched head rotations and body gestures that undermine viewer immersion. To address this limitation, we introduce sparse-frame video dubbing, a novel paradigm that preserves only reference keyframes while leveraging modern generative foundation models. As shown in \cref{fig:teaser}, it references only select keyframes to preserve the original video's emotional cadence, symbolic gestures, and camera trajectories, while liberating facial expressions, head motions, and body dynamics to synchronize organically with dubbed audio. As a long video generation task, it demands robust temporal continuation capabilities. A practical solution is employing audio-conditioned video generators with initial and terminal frame guidance.

% A practical implementation leverages audio-driven image-to-video models \cite{ltc, hallo3}, which typically generate short clips from a single initial reference frame. For long-form sequences, this method requires autoregressive chunk generation—where each new segment must initialize from the last generated frame of the preceding segment. As illustrated in \cref{fig:long_video_paradigm}, this gradual detachment from source material confines the entire sequence to just one reference frame, causing progressive deviation from the source video's identity and visual style. To counter this drift, we adapt models for initial/terminal frame conditioning, enabling explicit anchoring to multiple source keyframes $t_{\mathrm{ref}} > 1$ at strategic segment boundaries. This periodic realignment preserves the source's emotional expressions and gestures while accommodating extended generation.

% However, rigid enforcement of terminal frames ("hard conditioning") over-constrains motion trajectories. This forces generated movements to strictly interpolate between keyframes, suppressing audio-synchronized dynamics when source gestures conflict with speech rhythms. Sparse-frame dubbing thus necessitates soft conditioning—allowing flexible deviation from reference poses when required for audio alignment. The next subsection empirically evaluates these naive solutions, revealing their fundamental limitations and motivating our proposed approach.


% \paragraph
\subsection{Observation on naive solutions}
\begin{figure*}[t]
  \centering
  \includegraphics[width=1.0\textwidth]{res/observation1.pdf}
  \caption{
    (\textbf{left}): I2V model accumulates error for long video sequences. (\textbf{right}): A new chunk starts from frame 82. FL2V model suffers from abrupt inter-chunk transitions. 
  }
  \label{fig:observation1}
\end{figure*}
\begin{figure*}[t]
  \centering
  \includegraphics[width=1.0\textwidth]{res/observation2.pdf}
  \caption{
    % Analysis on how the similarity between the context frames and the reference frame influence the motion resemblance between the input and output video. The context frames for all the generated videos are shown in the left-bottom corner.
    A visual comparison between the training reference positioning strategies. All video chunks are generated using the same context frames and the same reference frame shown in below. 
  }
  \label{fig:observation2}
\end{figure*}
% \paragraph{Naive solutions to sparse-frame video dubbing}
% This section empirically evaluates practical approaches for long-form sparse-frame dubbing, revealing critical trade-offs between identity stability and motion naturalness. We begin with audio-driven image-to-video models \cite{let, cui2024hallo3}, which generate short clips (1-3 seconds) from a single initial reference frame. For extended sequences (10+ seconds), these require autoregressive chunk generation—splitting the video into segments (2-5 seconds each) where each new chunk initializes from the last generated frame of its predecessor. As \cref{fig:observation1} illustrates, this confines the entire sequence to one source reference. Consequently, \cref{fig:observation1} shows progressive degradation: facial identity drifts from the source actor, while backgrounds develop artifacts like color shifts or blurring due to lost visual details after the first chunk.

% To address identity drift, we adapt models for initial/terminal frame conditioning, anchoring each chunk to two source keyframes (start and end) as depicted in \cref{fig:observation1}. While this preserves expressions and gestures (\cref{fig:observation1}), it introduces a new limitation: when boundary frames appear similar—common in static scenes—motions become unnaturally subdued. Head rotations shrink to minimal arcs, gestures lose amplitude, and body language ignores audio energy, forcing rigid interpolation between poses. Crucially, both approaches exhibit abrupt jumps between chunks (\cref{fig:observation1}), as boundary-only conditioning ignores motion momentum—the kinetic continuity (e.g., a sustained head turn) that should flow across segments. In summary, these results expose a core dilemma: initial-frame conditioning preserves motion freedom but loses identity; dual-frame conditioning stabilizes identity but sacrifices expressiveness; both fail to achieve seamless transitions. This underscores the need for architectures that simultaneously maintain reference fidelity, motion naturalness, and temporal coherence.

This section investigates practical approaches for sparse-frame video dubbing using two baseline models: image-to-video (I2V) \cite{cui2024hallo3} and first-last-frame-to-video (FL2V) \cite{wan}. As illustrated in \cref{fig:observation1}, both methods exhibit critical limitations when generating long video sequences. The I2V approach operates by initializing the first chunk of the video from a single reference frame (e.g., the source video’s starting keyframe). For subsequent chunks, it uses only the last generated frame of the preceding chunk as the new reference. While this preserves motion flexibility, the lack of persistent anchoring to the original keyframes leads to accumulated errors: subtle discrepancies in identity (e.g., facial features gradually deviating from the source actor) and color tones (e.g., background hues shifting across chunks) compound over time, resulting in visible degradation. In contrast, the FL2V method conditions each chunk on both the start and end frames of the input segment, ensuring alignment with the source video’s reference poses. This eliminates accumulation errors but introduces a new problem: the model enforces rigid control by strictly replicating the reference frames at the corresponding timestamp. This contradicts the soft conditioning required for sparse-frame dubbing, where full body motions must dynamically adapt to audio cues.

Crucially, both methods suffer from abrupt inter-chunk transitions. Since they rely solely on static image conditions (e.g., a single frame for I2V or two fixed frames for FL2V), they lack momentum information that should carry over between chunks. These observations a highlight fundamental trade-off: I2V prioritizes motion fluidity at the expense of accumulated error, while FL2V prioritizes reference fidelity at the cost of motion naturalness.

% \paragraph{Analysis on reference-context similarity}
% To resolve abrupt transitions between video segments—jarring visual jumps caused by lost motion continuity—we implement a streaming generator that conditions on context frames from prior segments. This preserves motion momentum: the natural kinematic flow (e.g., a head turn mid-rotation) that should carry smoothly across chunks (implementation details in \cref{sec:pipeline}). Given the trade-offs observed in initial/terminal conditioning, we hypothesize that effective sparse-frame dubbing hinges on context-reference similarity—how similar it is between the reference frame and the context frames. 

% We therefore train a model accepting reference frames from arbitrary timestamps. Since the first chunk exhibits no accumulated error, it always uses the source’s first frame for stable initialization. For subsequent chunks, we evaluate four positioning strategies relative to the chunk start (illustrated in \cref{fig:observation2}):
% \begin{itemize}
%     \item Strategy 1 (start frame) imposes strong constraints, forcing exact reproduction of the reference at the video start. While stabilizing identity, it causes abrupt intra-chunk transitions because prior segments cannot perfectly match the new start frame.
%     \item Strategy 2 (excessively distant frames) induces out-of-domain artifacts, failing to generate coherent video as temporal gaps exceed the model’s comprehension.
%     \item Strategy 3 (terminal-proximal frames) applies intense conditioning at the reference timestamp, driving the model to reproduce its content near the ending frame. This alleviates accumulated error and resolves transitions (sufficient temporal gap enables smooth flow from transient frames), but triggers minimal motions by over-prioritizing pose matching.
%     \item Strategy 4 (start-proximal frames) positions references at moderate distances (within 2s of start), exerting balanced conditioning strength. The model avoids exact content reproduction at the reference timestamp, instead using it as a stabilizing guide. This achieves robust identity preservation, plausible motion, and seamless transitions.
% \end{itemize}

% As \cref{fig:observation2} confirms, Strategy 4’s success stems from correct reference frame proximity to the starting frame. It provides "corrective signals" for accumulated errors without stifling kinematic evolution from transient frames to the reference frame. 
% The inductive bias of training a random reference frame-conditioned streaming generator creates excessive constraint when references are positioned at starting/near-ending locations. This bias forces the model toward rigid reproduction of reference content, disrupting natural motion dynamics. Strategy 1 fails because forcing exact reproduction of the start frame creates rigid demands that disrupt motion continuity. Strategy 3 fails due to intense conditioning near endpoints, it generates minimal "safe" movements to match the reference frame that suppress natural gestures. Our expriments show the positioning of the reference image determines the behavior of the model. These findings provide valuable insights to build effective sparse-frame dubbing method by using fine-grained reference frame positioning in the long video sampling strategy.


% This subsection empirically evaluates naive approaches for sparse-frame video dubbing, revealing fundamental limitations in sparse-frame video dubbing.
% Audio-driven models using only initial frame conditioning exhibit progressive identity degradation and background distortion during extended sequences, as shown in \cref{fig:observation_accumulation}. This accumulation stems from referencing just one source frame at initialization, losing critical reference information for subsequent chunks discrete video segments typically spanning 2–5 seconds. As temporal distance increases, facial features diverge from source identity while backgrounds develop artifacts. Models incorporating initial/terminal frame conditioning mitigate identity drift through multi-keyframe anchoring but generate unnaturally subdued head and body motions when boundary frames exhibit high similarity—common in static dubbing scenes. These systems favor minimal transitions between keyframes, suppressing expressive gestures despite audio cues. Both approaches also produce abrupt temporal discontinuities between chunks (\cref{fig:observation_transition}), where single-frame conditioning fails to preserve motion continuity.

% To resolve discontinuities, we implement a streaming model (detailed in \cref{sec:streaming_model}) conditioning on transient frames from previous chunks. As demonstrated in \cref{fig:observation_streaming}, this maintains kinematic fluidity by inheriting motion context. To verify whether the minimal-motion is caused by the interpolating tendency between the initial and terminal frames. We remove the initial frame conditioning (omitting the initial reference constraint for chunks beyond the first segment) to allow motions to evolve organically rather than anchoring to fixed starting poses. This refinement eliminates the tendency toward subdued movements while retaining source identity through periodic keyframe anchoring. As quantified in \cref{fig:observation_strength}. Hard conditioning on the terminal frame restricts motion at certain frames, compromising motion-audio synchronization for sparse-frame video dubbing. To address this, we reduce terminal conditioning strength (detailed in \cref{sec:controling_conditioning_strength}) to enable dynamic deviation from reference poses during expressive speech segments. This flexibility improves lip and body synchronization without compromising identity stability. These insights collectively establish that effective sparse-frame dubbing requires streaming architectures with adaptive conditioning and motion-aware sampling to balance stability and expressiveness.

% Audio-driven models relying solely on initial frame conditioning exhibit progressive identity and background degradation during extended sequences (10–60 seconds), as visualized in \cref{fig:observation_accumulation}. Error accumulation occurs because these models reference only one source frame at initialization, losing critical reference information for subsequent chunks—discrete video segments typically spanning 2–5 seconds. This manifests as measurable facial distortion and color shifts that intensify with temporal distance from the starting frame.

% In contrast, approaches incorporating both initial and terminal frame conditioning mitigate identity drift through multi-frame anchoring. However, when initial and terminal frames exhibit high similarity—frequent in static dubbing scenes—they generate unnaturally minimal head and body motions. As quantified in \cref{fig:observation_strength}, these models default to rigid interpolation between keyframes, suppressing expressive gestures despite audio cues. Notably, both aforementioned solutions suffer abrupt transitions between chunks (\cref{fig:observation_transition}), where single-frame boundary conditioning fails to propagate motion momentum.

% Transition artifacts are resolved through streaming models with transient frame conditioning, which leverage sequences from prior chunks (typically 4–8 frames) to preserve kinematic continuity. As demonstrated in \cref{fig:observation_streaming}, this approach reveals a pivotal insight: skipping initial frame conditioning alleviates minimal-motion artifacts by eliminating interpolation constraints. Furthermore, reducing terminal conditioning strength—via techniques explored in \cref{sec:controlling_conditioning_strength}—enables flexible deviation from reference poses, significantly enhancing audio-motion synchronization while maintaining visual coherence.

% An analysis on why audio-driven image-to-video models fail to produce satisfing video dubbing results. We use audio-driven image-to-video models in two cases: (a) Using streaming audio-driven image-to-video model (initial frame conditioning). (b) Using streaming audio-driven image-to-video model with initial and terminal frame conditioning. 

%  The problems and reasons are: 
%  (1) Identity loss in (a). Phenomenon: The facial identity within the video shifts especially when the video frame has longer distance from the initial conditioning frame or the generated motion is large (e.g. from front face to side face). The video becomes oversaturated and the visual content deteriorates when the video becomes longer.  Cause: Existing image-to-video foundation models struggles to capture and maintain such identity; The reference frame is too sparse that some critical information is missing. Evidence: We test on our model, a figure shows the identity error and two possible fixes. When we apply initial and terminal frame conditioning, the problem is alleviated. When we decrease reference image interval, the problem is alleviated. 
%  (2) Stiff lip and body motion when the initial and terminal frames are similar in (b). Phenomenon: When the initial and terminal frames are similar (common in video dubbing, since the scene may not be changed) the lip and body motion is stiff and unnatural. Cause: the model tends to think the motion is limited and prone not produce the diverse and realistic motion. Evidence: As shown in the figure. We test our model, and Wan 2.1 (without audio condition driving, the situation of Wan 2.1 is worse), both exhibits the problem. But, when we remove one of the conditioning frame (either the inital or the terminal frame), the model will produce larger and realistic motion.

% \paragraph{Identity loss}
% Despite recent advances in audio-driven models enabling precise lip synchronization \cite{latentsync}, our investigation reveals fundamental limitations in streaming video dubbing implementations. When employing only initial frame conditioning, models exhibit progressive identity degradation during significant facial movements such as frontal-to-profile rotations. Identity accuracy systematically decreases as frames move farther from the reference, causing measurable face distortion and color shifts due to insufficient identity retention in foundation models and sparse reference frames missing critical viewpoint details. As demonstrated in \cref{fig:identity_error}, this limitation is mitigated through terminal frame augmentation or increased reference frequency, confirming that denser visual anchoring improves stability.

% \paragraph{Stiff motion}

% Building on this observation, we next examine systems using both initial and terminal frames—the common solution for temporal coherence. Unfortunately, these produce unnaturally rigid lip and body motions when frame similarity is high, as frequently occurs in static dubbing scenes. This motion impoverishment emerges because perceptual frame proximity triggers minimal-motion generation, suppressing expressive dynamics despite audio input. \cref{fig:motion_rigidity} documents this failure in our implementation and Wan 2.1 (where absent audio modulation exacerbates stiffness), while removing either conditioning frame restores natural motion diversity. These results expose a core architectural tension: while terminal frames reduce identity decay, they simultaneously constrain kinematic expressiveness when scenes lack visual variation.

% \paragraph{}
% \paragraph{Stiff motion when received similar initial and terminal frames}

\subsection{Audio-driven streaming video generator with reference frames}
% \paragraph{Building streaming video generator}
\label{sec:pipeline}
\begin{figure*}[t]
  \centering
  \includegraphics[width=1.0\textwidth]{res/infinitetalk.pdf}
  \caption{
    Visualization of InfiniteTalk pipeline. \textbf{Left}: The streaming model receives a audio, a reference frame, and context frames to denoise iteratively. \textbf{Right}: The architecture of the diffusion transformer. In addition to the traditional structures, each block includes an audio cross-attention layer and a reference cross-attention layer 
  }
  \label{fig:pipeline}
\end{figure*}
\begin{figure*}[t]
  \centering
  \includegraphics[width=1.0\textwidth]{res/strategies.pdf}
  \caption{
  Visualization of reference frame conditioning strategies for video dubbing models. Top four rows: conditioning on input video frames. Bottom row: conditioning on generated video frames. \textbf{Left}: Image-to-video dubbing model with initial frame conditioning (I2V) and initial+terminal frame conditioning (IT2V). \textbf{Right}: Streaming dubbing model with four conditioning strategies. Within each category (left/right), all strategies share identical generated-video conditioning approaches.
    % A sketch for reference frame conditioning for image-to-video dubbing model and our streaming model. In the first four rows, we show which frame is conditioned on the input video. In the last row, the conditioning on the generated video is illustrated. \textbf{(left)} shows both initial frame conditioning (I2V) and initial & terminal frame conditioning (IT2V). They use the same generated video conditioning strategy. \textbf{(right)} shows the four conditioning strategies on the streaming dubbing model. They share the same generated video conditioning strategy. 
  }
  \label{fig:mps}
\end{figure*}
% Building a streaming audio-driven image-to-video diffusion. We build a audio -driven image-to-video generator model. The model consists: video DiT (each transformer block includes self-attention, audio cross-attention, and reference image cross-attention). Then introduce the detailed streaming context and reference image injection method.
To resolve the accumulated accumulated error in I2V models and the abrupt transitions in FL2V models, we build a audio-driven streaming human animation architecture. This framework employs context frames, defined as the trailing segment of each previously generated chunk, to propagate kinetic momentum into subsequent segments. By processing these frames through a diffusion transformer, the model sustains motion continuity. To eliminate accumulated errors, we adopt multiple reference frames dynamically sampled from the source video, similar to FL2V's multi-frame conditioning. These keyframes preserve critical visual attributes including identity, background details, camera trajectories, and stylistic elements. Crucially, unlike FL2V's rigid replication of reference frames at fixed positions, which suppresses natural motion. Our model has the potential to achieve soft conditioning as discussed in \cref{sec:strategies}.
% We develop an audio-driven human animation that has the capability to perform streaming long video generation and reference frame conditioning. 

The sketch for the model is shown in \cref{fig:pipeline}. The model consists an audio embedder \cite{wav2vec}, a video VAE, and diffusion transformer (DiT) \cite{dit}. We first introduce how our model is trained. In the followings, let the video latent refers to embedding derived by encoding pixel space video by using the video VAE. Similarly, let audio embedding denote the embedding computed by encoding a audio sequence using the audio embedder. To train the model, we do not need a dubbed video pair, but only a video with the audio track is needed. Given a source video. The reference frame is uniformly random sampled from the video. During training, the context frames are the first $4(t_c - 1) + 1$ frames of the source video. After VAE encoding, we derive reference frame latent $\boldsymbol{x}_{\mathrm{ref}} \in \mathbb{R}^{c \times 1 \times h \times w}$, the full source video latent $\boldsymbol{x}_{\mathrm{full}} \in \mathbb{R}^{c \times (t + t_c) \times h \times w}$, the context frames latent $\boldsymbol{x}_{\mathrm{context}} \in \mathbb{R}^{c \times t_c \times h \times w}$, and the subsequent frames latent $\boldsymbol{x}_0 \in \mathbb{R}^{c \times t \times h \times s}$ separately, where the full video latent is a combination of the context frames latent  $\boldsymbol{x}_{\mathrm{context}}$ and the subsequent frames latent $\boldsymbol{x}_0$, $\boldsymbol{x}_{\mathrm{full}} = \{\boldsymbol{x}_{\mathrm{context}}, \boldsymbol{x}_0\}$. The audio sequence in the source video is encoded to get embedding $a$. Without a loss of generality, we show the process to train at $t$ in conditional flow matching in this section. Unit gaussian noise is used as the trivial distribution. The noisy latent is derived by $\boldsymbol{x}_t = (1 - t) \cdot \boldsymbol{x}_1 + t \cdot \boldsymbol{x}_0, \boldsymbol{x}_1 \sim \mathcal{N}(\boldsymbol{0}, \mathbb{I})$ where $\boldsymbol{x}_1$ is gaussian noise that has the same dimensionality of $\boldsymbol{x}_0$. The DiT model formulates a field estimator $\boldsymbol{v}_\theta(\boldsymbol{x}_t \vert \boldsymbol{c})$ using condition $\boldsymbol{c} = \{y, a, \boldsymbol{x}_{\mathrm{ref}}, \boldsymbol{x}_{\mathrm{tran}}\}$. Specifically, we first concatenate the noisy latent $\boldsymbol{x}_t$ and the clean context frames $\boldsymbol{x}_{\mathrm{context}}$ in the temporal dimension to get $\boldsymbol{z}_1 \in \mathbb{R}^{c \times (t + t_c) \times h \times w}$. Then we pad the reference frame to the temporal length $t_c + t$ to get $\boldsymbol{z}_2 \in \mathbb{R}^{c \times (t + t_c) \times h \times w}$. Finally, we concatenate $\boldsymbol{z}_1$, $\boldsymbol{z}_2$ and a reference frame indicating mask $\boldsymbol{m} \in \mathbb{R}^{4 \times (t + t_c) \times h \times w}$ in channel dimension. Mathematically, the process is

\begin{equation}
\begin{aligned}
    \boldsymbol{z}_1 & = \mathrm{concat}((\boldsymbol{x}_{\mathrm{context}}, \boldsymbol{x}_t), 2) \\
    \boldsymbol{z}_2 & = \mathrm{concat}((\boldsymbol{x}_{\mathrm{ref}}, \mathbb{0}), 2) \\
    \boldsymbol{m}   & = \mathrm{concat}((\mathbb{1}, \mathbb{0}), 2) \\
    \boldsymbol{z}   & = \mathrm{concat}((\boldsymbol{z}_1, \boldsymbol{z}_2, \boldsymbol{m}), 1)
\end{aligned}
\end{equation}

where $\mathrm{concat}(\cdot)$ is the concatenation operator (e.g, $\mathrm{concat}((\boldsymbol{x}_{\mathrm{context}}, \boldsymbol{x}_t), 2)$ concatenates $\boldsymbol{x}_{\mathrm{context}}$ and $\boldsymbol{x}_t$ in the 2-nd dimension). $\mathbb{0}$ and $\mathbb{1}$ are zero and one tensors that have the dimensions $\mathbb{0} \in \mathbb{R}^{4 \times (t + t_c - 1) \times h \times w}, \mathbb{1} \in \mathbb{R}^{4 \times 1 \times h \times w}$. As depicted by \cref{fig:pipeline} \textbf{(right)}. Within our transformer model, there is an audio cross-attention module and an image cross-attention that achieves audio and reference image conditioning. The reference frame is processed by CLIP vision model \cite{openclip} to get the embedding $\boldsymbol{z}_{\mathrm{ref}}$ before feeding to the DiT. To train this model, we adopt the conditional flow matching objective \cite{flowmatching}.

\begin{equation}
    \mathcal{L}_{\mathrm{fm}} = \mathbb{E}_{t, \boldsymbol{x}_0, \boldsymbol{x}_1, \boldsymbol{c}} \|\boldsymbol{v}_\theta(\boldsymbol{x}_t \vert \boldsymbol{c}) - (\boldsymbol{x}_1 - \boldsymbol{x}_0)\|_2^2.
\end{equation}
% Given a source video, a audio sequence, a reference frame sampled randomly from the video, and the transient frames from the previous video segment. We use the audio embedder, and the video VAE to encode them. 

% We also propose a weak reference model that reduces the dependency of the velocity output from the reference frame. 
% During training, the weak reference model samples reference frame from instead of from the current source video chunk.

Then, we briefly introduce the sampling method. An illustration is presented in \cref{fig:mps}. We generate the whole long video sequence by auto-regressively generating small video chunks. In the first video chunk, we use the first frame of the input video as the reference frame and no context frames are required. In the following video chunks, we use the last $4 (t_c - 1) + 1$ frames from the previous output video chunk as the context frames and use the first image of the input chunk as the reference frame.
% We design four sampling strategies of reference frames to achieve a pleasant trade-off between control strength and synchronization.

\subsection{Soft conditioning}
\label{sec:strategies}
\paragraph{Control strength from the Reference frame}
This section explores strategies for achieving soft conditioning in sparse-frame video dubbing where the model must generate audio-aligned full-body motions without rigidly replicating reference frames that may conflict with dubbed speech. Meanwhile, we expect the model to have adaptive control strength: (1) when the reference is similar to the context frames, the control strength is weak such that the model produces diverse dynamics. (2) when the reference is very distinct to the context frames, the control strength is stronger to ensure better consistency in identity and background. Below, we show the analysis on the training strategies that match the requirements.

We initiate our investigation with Model M0, which samples reference frames uniformly at random from the current source input video chunk during training. As demonstrated in \cref{fig:observation2}, this approach exhibits excessive control strength. The model inappropriately duplicates reference content at arbitrary timestamps, such as replicating slapping on the forehead during emotionally neutral speech, disrupting audio-visual synchronization.

To systematically analyze how reference positioning during training governs control fidelity, we examine two granular dimensions: chunk-level positioning (selecting which temporal segment provides the reference) and frame-level positioning (choosing the specific frame within that segment). We train three model variants to isolate these effects. Model M1 samples references exclusively from the first or last frame of the input chunk. This strategy mirrors FL2V's rigidity, forcing the generated video to replicate reference poses precisely at chunk boundaries, even when they contradict the audio's emotional cadence. Model M2 samples references from temporally distant chunks (e.g., segments separated by >5 seconds). While this weakens control sufficiently to avoid replication, it introduces accumulated color and background errors over long sequences, indicating insufficient fidelity preservation. In contrast, Model M3 samples references from adjacent chunks (e.g., within 1 seconds of the input). This configuration achieves moderate control strength: references preserve identity and camera motion without exact duplication, while eliminating accumulated errors entirely.

Our experiments conclusively demonstrate that chunk-level distance is the dominant factor modulating control strength. Shorter temporal distances (as in M3) create an optimal equilibrium: they anchor visual consistency to the source while liberating facial expressions, head rotations, and body gestures to organically synchronize with audio. Longer distances (as in M2) degrade preservation capabilities, leading to destabilized outputs. Fixed boundary sampling (as in M1) prioritizes replication over expressiveness, stifling motion dynamics. Thus, M3’s near-chunk positioning emerges as the foundational strategy for soft conditioning—enabling faithful yet flexible video dubbing where motions breathe in harmony with speech. 
% Our experiments reveal that chunk-level distance is the primary modulator of control strength. Shorter temporal distances (as in M3) yield balanced conditioning—preserving source attributes while enabling natural, audio-driven motion—whereas longer distances (M2) degrade preservation, and boundary fixation (M1) stifles expressiveness. M3’s near-chunk positioning thus emerges as the optimal strategy for soft conditioning in sparse-frame dubbing.

% Based on our observation and analysis. We propose a fine-grained sampling strategy to ensure satisfying accumulated error correction and audio synchronized lip, head, and body motion generation. An illustration is presented in \cref{fig:mps}. We generate the whole long video sequence by auto-regressively generating small video chunks. In the first video chunk, we use the first frame of the input video as the reference frame and no context frames are required. In the following video chunks, we use the last $4 (t_c - 1) + 1$ frames from the previous output video chunk as the context frames. We design four sampling strategies of reference frames to achieve a pleasant trade-off between control strength and synchronization.

% \begin{itemize}
%     \item Strategy 1 (first frame): Using the first frame of the input video chunk as the reference frame. The context-reference similarity will be the highest. This strategy aims to provide the weakest control from the input video.
%     % imposes strong constraints, forcing exact reproduction of the reference at the video start. While stabilizing identity, it causes abrupt intra-chunk transitions because prior segments cannot perfectly match the new start frame.
%     \item Strategy 2 (last frame): Using the last frame as the reference frame. The expected context-reference similarity is low, providing a strong control over the motion.
%     % induces out-of-domain artifacts, failing to generate coherent video as temporal gaps exceed the model’s comprehension.
%     \item Strategy 3 (far frame): The weak reference model is used in this strategy. The reference frame is picked from a timestamp far from the current video chunk. It provides minimal control over the video motion but preserves human identity.
%     % applies intense conditioning at the reference timestamp, driving the model to reproduce its content near the ending frame. This alleviates accumulated error and resolves transitions (sufficient temporal gap enables smooth flow from transient frames), but triggers minimal motions by over-prioritizing pose matching.
%     \item Strategy 4 (proximal first frame): Using the frame near the first frame (within 1 second) as the reference frame. It is likely to create a medium reference-context similarity, providing a moderate motion resemblance between the input and output video.
%     % positions references at moderate distances (within 2s of start), exerting balanced conditioning strength. The model avoids exact content reproduction at the reference timestamp, instead using it as a stabilizing guide. This achieves robust identity preservation, plausible motion, and seamless transitions.
% \end{itemize}


% Our experiments show strategy 4 achieves the best performance in sparse-frame video dubbing. It achieves a moderate motion control strength. Ensuring plausible lip, body, and expression synchronization while maintaining identity, iconic gestures, backgrounds, and camera movements. 

% We achieve a moderate control strength from the reference frame by 

% \paragraph{Camera trajectory injection?}
\paragraph{Camera control}
% \subsection{Camera Control}
We explore ways to address the camera movement preservation in sparse-frame video dubbing in this section. The utilization of reference frames is providing a global control of camera trajectory. However, the detailed camera movement within each video chunk is not controlled and may contradict the source video. To resolve this, we use two plugin including SDEdit \cite{sdedit} and Uni3C \cite{cao2025uni3c}. SDEdit incorporates trajectory information by adding the source video to the initialize noise at a scale $t_0$. The denoising sampling process starts from $t= t_0$ instead of $t=1$. $\boldsymbol{x}_{t_0} = (1 - t_0) \cdot \boldsymbol{x}_1 + t_0 \cdot \boldsymbol{x}_0$
Uni3C injects camera movements by deploying a ControlNet-like architecture. A comparison between the methods is shown in our experiments.
